\begin{lemma}
\label{lemma-coherent-filter}
Let $X$ be a Noetherian scheme.
Let $\mathcal{F}$ be a coherent sheaf on $X$.
There exists a filtration
$$
0 = \mathcal{F}_0 \subset \mathcal{F}_1 \subset
\ldots \subset \mathcal{F}_m = \mathcal{F}
$$
by coherent subsheaves such that for each $j = 1, \ldots, m$
there exist an integral closed subscheme $Z_j \subset X$ and
a nonzero coherent sheaf of ideals $\mathcal{I}_j \subset \mathcal{O}_{Z_j}$
such that
$$
\mathcal{F}_j/\mathcal{F}_{j - 1}
\cong (Z_j \to X)_* \mathcal{I}_j
$$
\end{lemma}

\begin{proof}
Consider the collection
$$
\mathcal{T} =
\left\{
\begin{matrix}
Z \subset X
\text{ closed such that there exists a coherent sheaf }
\mathcal{F} \\
\text{ with }
\text{Supp}(\mathcal{F}) = Z
\text{ for which the lemma is wrong}
\end{matrix}
\right\}
$$
We are trying to show that $\mathcal{T}$ is empty. If not, then
because $X$ is Noetherian we can choose a minimal element
$Z \in \mathcal{T}$. This means that there exists a coherent
sheaf $\mathcal{F}$ on $X$ whose support is $Z$ and for which the
lemma does not hold. Clearly $Z \not = \emptyset$ since the only
sheaf whose support is empty is the zero sheaf for which the
lemma does hold (with $m = 0$).

\medskip\noindent
If $Z$ is not irreducible, then we can write $Z = Z_1 \cup Z_2$
with $Z_1, Z_2$ closed and strictly smaller than $Z$.
Then we can apply Lemma \ref{lemma-prepare-filter-support}
to get a short exact sequence of coherent sheaves
$$
0 \to
\mathcal{G}_1 \to
\mathcal{F} \to
\mathcal{G}_2 \to 0
$$
with $\text{Supp}(\mathcal{G}_i) \subset Z_i$. By minimality of
$Z$ each of $\mathcal{G}_i$ has a filtration as in the statement
of the lemma. By considering the induced filtration on $\mathcal{F}$
we arrive at a contradiction. Hence we conclude
that $Z$ is irreducible.

\medskip\noindent
Suppose $Z$ is irreducible. Let $\mathcal{J}$ be the sheaf of ideals
cutting out the reduced induced closed subscheme structure of $Z$,
see Schemes, Lemma \ref{schemes-lemma-reduced-closed-subscheme}.
By Lemma \ref{lemma-power-ideal-kills-sheaf} we see there exists
an $n \geq 0$ such that $\mathcal{J}^n\mathcal{F} = 0$. Hence we obtain
a filtration
$$
0 = \mathcal{J}^n\mathcal{F} \subset \mathcal{J}^{n - 1}\mathcal{F}
\subset \ldots \subset \mathcal{J}\mathcal{F} \subset \mathcal{F}
$$
each of whose successive subquotients is annihilated by $\mathcal{J}$.
Hence if each of these subquotients has a filtration as in the statement
of the lemma then also $\mathcal{F}$ does. In other words we may
assume that $\mathcal{J}$ does annihilate $\mathcal{F}$.

\medskip\noindent
In the case where $Z$ is irreducible and $\mathcal{J}\mathcal{F} = 0$
we can apply Lemma \ref{lemma-prepare-filter-irreducible}.
This gives a short exact sequence
$$
0 \to
i_*(\mathcal{I}^{\oplus r}) \to
\mathcal{F} \to
\mathcal{Q} \to 0
$$
where $\mathcal{Q}$ is defined as the quotient.
Since $\mathcal{Q}$ is zero in a neighbourhood of $\xi$ by
the lemma just cited we see that the support of $\mathcal{Q}$
is strictly smaller than $Z$. Hence we see that $\mathcal{Q}$
has a filtration of the desired type by minimality of $Z$.
But then clearly $\mathcal{F}$ does too, which is our final contradiction.
\end{proof}